\documentclass[10pt,a4paper]{amsart}
\usepackage[margin=19mm]{geometry}
\usepackage{amsmath,amssymb,mathtools}
\usepackage[scr=rsfs,cal=euler]{mathalfa}
\usepackage{xfrac}
\usepackage[unicode,hidelinks,bookmarks=false]{hyperref}
\newcommand{\ie}{i.e.\ }
\newcommand{\cf}{cf.\ }
\newcommand{\resp}{respectively\ }
\newcommand{\Subsection}{\subsection}
\providecommand{\nobreakdash}{\nobreak\hbox{-}}
\setlength{\emergencystretch}{2em}
\multlinegap0pt
\usepackage{bm}
\usepackage{bbm}
\usepackage{dsfont}
\usepackage{xspace}
\usepackage{cancel}
\usepackage{mathtools}
\usepackage{cleveref}
\usepackage{amsthm}
\makeatletter
\@ifundefined{proposition}{\newtheorem{proposition}{Proposition}}{}
\makeatother
\makeatletter
\newcommand{\R}{\mathbb{R}}    % Real numbers
\newcommand{\Star}[2]{{{#1}^{#2}}}
\newcommand{\arroid}{\mathbf{A}}
\newcommand{\arroidDeletion}[1]{\arroid\setminus #1}
\DeclareMathOperator{\cone}{cone}
\newcommand{\groundset}{\mathcal{A}}
\expandafter\ifx\csname introthm\endcsname\relax
\newenvironment{introthm}[1]{\stepcounter{introthmnum}
    \newtheorem*{introthmTemp:\arabic{introthmnum}}{#1}\begin{introthmTemp:\arabic{introthmnum}}}{\end{introthmTemp:\arabic{introthmnum}}}
\fi
\newcommand{\mulSet}{m}
\newcommand{\pbf}{\mathbf{p}}
\newcommand{\points}{\mathcal{P}}
\DeclarePairedDelimiter{\set}{\lbrace}{\rbrace}        % Set
\DeclareMathOperator{\tm}{TM}
\makeatother
\title[Tropicalization of curve arrangement complements and arroids]{Tropicalization of curve arrangement complements and arroids}
\author{Edvard Aksnes}
\date{Source: arXiv:2404.14380v3 (CC BY 4.0)}
\begin{document}
\maketitle

\noindent\emph{Source and licence.} Tropicalization of curve arrangement complements and arroids. Version arXiv:2404.14380v3 (CC BY 4.0), distributed under the Creative Commons Attribution 4.0 International licence, as recorded in the arXiv licence metadata for this exact version and shown on the abstract page https://arxiv.org/abs/2404.14380v3. Full rights notice: licenses/arxiv-2404.14380-CC-BY-4.0.txt.\par\smallskip

\noindent\textit{Editorial scope.} Counted unit: the Proposition whose source label is \texttt{prop:arroid\_tropical\_modif} in the article (source0.tex lines 810--812, source label \texttt{prop:arroid\_tropical\_modif}), delivered together with the article's own complete original proof of it (source0.tex lines 813--833, one \texttt{proof} environment of 21 lines). No mathematics is written, repaired or re-derived: statement and proof are the article's own text line for line. Required context retained uncounted: prop:arroid\_star (lines 801--803). Every definition, display and label of those lines is retained unchanged. Omitted from this package and not redistributed: 1--800, 804--809, 834--1196. Nothing inside a retained excerpt is omitted or altered: every retained body is delivered exactly as the article's lines are written, so no omission receipt of any kind accompanies this package. Citations: the retained excerpts carry no citation command at all, so no bibliography entry of the article is transcribed and no reference is redistributed. Reference adaptation: none was needed and none was made; every reference inside the delivered unit resolves inside it, so no unresolved reference or citation is produced. Printed slips kept literal: no printed word, symbol, hypothesis or number of the counted statement or of its proof was altered. Layout and class: the article's own class is \texttt{mathscan} with packages including inputenc, fontenc, lmodern, microtype, amssymb, amsthm, thmtools, mathtools, mathrsfs, xcolor, graphicx, tikz, tikz-cd, tikz-3dplot; the wrapper is the lane's pinned \texttt{amsart} editorial wrapper at 10pt on A4, which restates the theorem-like environments the retained text needs and adds the article's own macro definitions that the retained text needs (\texttt{\textbackslash R}, \texttt{\textbackslash Star}, \texttt{\textbackslash arroid}, \texttt{\textbackslash arroidDeletion}, \texttt{\textbackslash cone}, \texttt{\textbackslash groundset}, \texttt{\textbackslash mulSet}, \texttt{\textbackslash pbf}, \texttt{\textbackslash points}, \texttt{\textbackslash set}, \texttt{\textbackslash tm}), with extra wrapper packages (bm, bbm, dsfont, xspace, cancel, mathtools, cleveref). The delivered numbering is the unit's own. Assets: no retained span contains a figure environment, a graphics-inclusion command, a PDF-page inclusion, a table or a TikZ picture; no image file is copied into this package, the package declares no asset dependency and no image file is redistributed with it. Licence: version arXiv:2404.14380v3 of this article is distributed under the Creative Commons Attribution 4.0 International licence, as recorded in the arXiv licence metadata for this exact version and shown on the abstract page https://arxiv.org/abs/2404.14380v3. A full rights notice is delivered at licenses/arxiv-2404.14380-CC-BY-4.0.txt. Characters: every retained excerpt is pure ASCII, so the delivered engine, XeLaTeX with its default Latin Modern text font, reports no missing character.
\begin{proposition}\label{prop:arroid_star}
    Let $\Sigma_\arroid$ be an arroid fan, and $i\in \groundset$ correspond to the ray $\rho_i$ of $\Sigma_\arroid$. Then the reduced star $\Star{\Sigma_\arroid}{\rho_i}$ is equal to the arroid fan $\Sigma_{\arroid / i}$ of the arroid contraction $\arroid / i$.
\end{proposition}

\begin{proposition}\label{prop:arroid_tropical_modif}
    Let $\arroid=(\groundset, d, \points, \mulSet)$ be a transversal arroid, and $\arroidDeletion{i}$ one of its arroid deletions. Then $\Sigma_\arroid = \tm(\Sigma_{\arroidDeletion{i}}, \Sigma_{\arroid / i})$ is a tropical modification of the fan of the arroid deletion $\arroidDeletion{i}$ along the fan of the arroid contraction $\arroid / i$.
\end{proposition}

\begin{proof}
    We will show that the preimage of a point in the fan $\Sigma_{\arroid}$ under the coordinate projection $\pi \colon \R^{|\groundset|}/\langle(d_1,\dots, d_n)\rangle \to \R^{|\groundset\setminus i|}/\langle(d_1,\dots, \hat{d_i} ,\dots ,d_n) \rangle$ is either a half-line in the $e_i$ direction, or a point. Then one may obtain a piecewise integer affine function $g$ defining the tropical modification by taking $g(p) = \pi^{-1}(p)$ for the points $p\in \Sigma_{\arroidDeletion{i}}$ whose preimage $\pi^{-1}(p)$ are single points. 

    To show that $\pi^{-1}(p)$ is either a point or half-line in $\Sigma_{\arroid}$, we analyze the cones of the latter. The cones of $\Sigma_{\arroid}$ have one of three forms: 
    \begin{itemize}
        \item $\cone(e_j, e_j+e_i)$ for some $j\in \groundset$, corresponding to a point $\pbf$ which disappears in $\arroidDeletion{i}$,
        \item $\cone(e_j, e_\mathbf{a} + e_i)$, corresponding to a point $\pbf = \mathbf{a}\cup {i} \in \points$ from which $i$ is removed in $\arroidDeletion{i}$, or 
        \item $\cone(e_j, e_\pbf)$ for some $\pbf\in \points$ not containing $i$.
    \end{itemize}

    Any point $p$ of the fan $\Sigma_{\arroidDeletion{i}}$ is contained inside $\cone(e_k,e_\pbf)$ for some $k$ and $\pbf\in \points \setminus i$ with $k \in \pbf$. We distinguish multiple cases:
    \begin{itemize}
        \item If $p=\mathbf{0}$ is the origin, then $\pi^{-1}(p)=\cone(e_i)$ is a half-ray of the fan $\Sigma_{\arroid}$.
        \item If $p = \alpha e_k + \beta e_\pbf$ with $\alpha,\beta >0$, and $\pbf \cup \set{i}$ is a point of $\arroid$, the preimage $\pi^{-1}(p)$ is $p' = \alpha e_j + \beta (e_\pbf+ e_i)$ in $\Sigma_{\arroid}$.
        \item If $p= \alpha e_k$ then either $\set{i,k}\in \points$, and then $\pi^{-1}(p)= \set{a e_j +b (e_j+e_i) \mid a+b=\alpha; \; a,b \geq 0}$ is a half-ray or $\set{i,k} \not\in \points$ and then $\pi^{-1}(p)=\alpha e_k$.
        \item If $p = \beta e_\pbf$ then either $\pbf\cup \set{i}$ is a point of $\arroid$, and the preimage is a half.line, or it is not so that $\pi^{-1}(p)=\beta e_\pbf$.
    \end{itemize}
    For any of $p$, the fiber is either a half-line or a point, and thus $\Sigma_{\arroid}$ is a tropical modification of $\Sigma_{\arroid \setminus i}$.
    
    Moreover, it follows from the construction that the divisor along which the modification is performed is the reduced star $\Star{\Sigma_\arroid}{\rho_i}$, which is $\Sigma_{\arroid/ i}$ by \cref{prop:arroid_star}, so that we have $\Sigma_\arroid = \tm(\Sigma_{\arroidDeletion{i}}, \Sigma_{\arroid / i})$. Moreover, the weights of $\Sigma_\arroid$ are then given in terms of the weights of the rays of $\Sigma_{\arroid / i}$, which is compatible with the construction of both fans as above.
\end{proof}

\end{document}
